\indent ವುದು. ಅಂದರೆ, ಕ್ರಿಯಾಪದವು ಭವಿಷ್ಯತ್‌ ಕಾಲದಲ್ಲಿರುವುದೆಂಬುದನ್ನು ತಿಳಿಸುವುದು.\par

\indent ಮೇಲಿನ ವಾಕ್ಯಗಳಲ್ಲಿ ಕಂಡುಬರುವ ಕ್ರಿಯಾಪದಗಳಿಂದ ಇವು ಮುಂದೆ ನಡೆಯುವ ಕ್ರಿಯೆಯನ್ನು ತಿಳಿಸುವ ಕ್ರಿಯಾಪದಗಳೆಂದು ತಿಳಿದುಬರುವುದು ಆದುದರಿಂದ ಈ ಕ್ರಿಯಾಪದಗಳಿಗೆ ಭವಿಷ್ಯತ್‌ ಕಾಲದ ಕ್ರಿಯಾಪದಗಳೆಂದು ಹೆಸರು. ಅದರಂತೆಯೇ ಮೇಲಿನ ಕ್ರಿಯಾಪದಗಳಲ್ಲಿರುವ ಅನು, ಅಳು, ಉದು, ಎ, ಎನು, ಅರು, ಉವು, ಇರಿ, ಎವು—ಎಂಬುವು ಭವಿಷ್ಯತ್‌ ಕಾಲದ ಆಖ್ಯಾತ ಪ್ರತ್ಯಯಗಳು.\par

\vspace{2.5pt plus 3pt minus 0.5pt}\noindent\hspace*{1.2em}\kbvnum{(4)}\hspace{0.55em}ವಿಧ್ಯರ್ಥಕ ಕ್ರಿಯಾಪದಗಳು\par

\vspace{3pt plus 2pt}\noindent\hspace*{1.2em}\kbvlabel{ಉದಾ}\par\vspace{1pt}

\indent ಅವರು ಮನೆಗೆ ಬರಲಿ .... (ಬರು + ಅಲಿ)\par

\indent ಸೀತೆಯು ಶಾಲೆಗೆ ಹೋಗಲಿ .... (ಹೋಗು + ಅಲಿ)\par

\indent ನೀವು ಸುಖವಾಗಿ ಬಾಳಿರಿ .... (ಬಾಳು + ಇರಿ)\par

\indent ನಾನು ಊಟ ಮಾಡುವೆ .... (ಮಾಡು + ಎ)\par

\indent ನಾವು ಹಾಡು ಹೇಳೋಣ .... (ಹೇಳು + ಓಣ)\par

\indent ಮೇಲಿನ ಉದಾಹರಣೆಗಳಲ್ಲಿ ಧಾತುಗಳೊಂದಿಗೆ ಅಲಿ, ಇರಿ, ಎ, ಓಣ\par

\indent —ಈ ಆಖ್ಯಾತ ಪ್ರತ್ಯಯಗಳು ಸೇರಿ ಕ್ರಿಯಾಪದಗಳಾಗಿರುವುವು. ಅಪ್ಪಣೆ, ಆಜ್ಞೆ, ಆಶೀರ್ವಾದ, ಹಾರೈಕೆ, ಸಮ್ಮತಿ ಕೊಡುವ ಸಂದರ್ಭಗಳಲ್ಲಿ ಬಳಸುವ ಈ ಕ್ರಿಯಾಪದಗಳಿಗೆ ವಿಧ್ಯರ್ಥಕ ಕ್ರಿಯಾಪದಗಳೆಂದು ಹೆಸರು. ಈ ಕ್ರಿಯಾ ಪದಗಳಲ್ಲಿ ಕಾಲಸೂಚಕ ಪ್ರತ್ಯಯ ಸೇರಿರುವುದಿಲ್ಲ.\par

\vspace{2.5pt plus 3pt minus 0.5pt}\noindent\hspace*{1.2em}\kbvnum{(5)}\hspace{0.55em}ನಿಷೇಧಾರ್ಥಕ ಕ್ರಿಯಾಪದಗಳು\par

\noindent\hspace*{1.2em}\kbvlabel{ಉದಾ}\hspace{0.6em}“ಮಾಡು” ಧಾತು.\par

\vspace{5pt plus 4pt minus 1pt}
\begin{center}
{\renewcommand{\arraystretch}{1.24}\setlength{\tabcolsep}{10pt}\rowcolors{2}{kbvrow}{white}
\begin{tabularx}{0.88\textwidth}{X X}
{\kbvsecfont\color{kbvaccentdark}ಏಕವಚನ} & {\kbvsecfont\color{kbvaccentdark}ಬಹುವಚನ} \\[1pt]
\arrayrulecolor{kbvline}\specialrule{0.5pt}{0pt}{2pt}\arrayrulecolor{black}
ಮಾಡು + ಅನು = ಮಾಡನು & ಮಾಡು + ಅರು = ಮಾಡರು \\
ಮಾಡು + ಅಳು = ಮಾಡಳು & ಮಾಡು + ಅರು = ಮಾಡರು \\
\end{tabularx}}
\end{center}
\vspace{5pt plus 4pt minus 1pt}
